\chapter{Joint Distributions, Covariance, and Correlation}
\label{ch:joint}

So far we have studied one random variable at a time. But many questions
concern \emph{two or more} quantities together: hair color and eye color, a
chip's defect status across a batch, a wood's density and its hardness. The
\emph{joint distribution} encodes how a pair of random variables behave
simultaneously, and from it we can read off whether they move together, and how
strongly.

\section{Joint and Marginal Distributions}

For discrete $X$ and $Y$, the \emph{joint pmf} is $p(a,b)=\Prob(X=a,Y=b)$. From
it we recover each variable's own (\emph{marginal}) distribution by summing out
the other:
\[
  p_X(a)=\sum_b p(a,b),\qquad p_Y(b)=\sum_a p(a,b).
\]
The marginals live in the row and column totals of the joint table.

\begin{example}[Completing a joint table]
The joint distribution of $X\in\{0,1,2\}$ and $Y\in\{-1,1\}$ is partly given:
the margins are $p_X=(\tfrac16,\tfrac23,\tfrac16)$ and
$p_Y=(\tfrac12,\tfrac12)$, the row $Y=1$ reads $(0,\tfrac12,0)$. Complete the
table and decide whether $X$ and $Y$ are independent.

\solution Each column must sum to its $X$-margin, so the row $Y=-1$ is
$(\tfrac16,\tfrac16,\tfrac16)$:
\[
\begin{array}{c|ccc|c}
 & 0 & 1 & 2 & p_Y\\\hline
 Y=-1 & 1/6 & 1/6 & 1/6 & 1/2\\
 Y=1 & 0 & 1/2 & 0 & 1/2\\\hline
 p_X & 1/6 & 2/3 & 1/6 & 1
\end{array}
\]
They are \emph{dependent}: for instance
$\Prob(X=0,Y=1)=0\neq\tfrac16\cdot\tfrac12=\Prob(X=0)\Prob(Y=1)$.
\end{example}

\begin{example}[Reconstructing a joint distribution]
$X$ and $Y$ each take values $1,\dots,5$ with margins
$p_X=(\tfrac1{14},\tfrac5{14},\tfrac4{14},\tfrac2{14},\tfrac2{14})$ and
$p_Y=(\tfrac5{14},\tfrac4{14},\tfrac2{14},\tfrac2{14},\tfrac1{14})$. Each joint
probability is either $0$ or $\tfrac1{14}$. Find the joint distribution.

\solution Each cell is $\tfrac1{14}$ or $0$, and the number of $\tfrac1{14}$
cells in a row (column) must equal $14$ times that row's (column's) margin.
Working row by row—five entries in the row summing to $\tfrac5{14}$, four in the
next, and so on, while respecting the column totals—forces a unique
``staircase'' pattern in which row $i$ has $14p_Y(i)$ filled cells packed toward
the columns with the largest remaining demand. The filled cells are exactly
those $(a,b)$ for which the cumulative column and row counts still admit a unit,
producing the table whose row sums are $5,4,2,2,1$ (times $\tfrac1{14}$) and
column sums $1,5,4,2,2$.
\end{example}

\begin{example}[Hair color and eye color]
The hair and eye color of $5383$ people are cross-classified. Dividing the
counts by $5383$ turns the table into a joint distribution for hair color $X$
and eye color $Y$. Are $X$ and $Y$ independent?

\solution The joint and marginal probabilities are the cell counts and
row/column totals divided by $5383$. Independence fails dramatically: the
observed $\Prob(\text{fair hair, light eyes})=1168/5383\approx 0.22$, but the
product of margins is $(1741/5383)(2298/5383)\approx 0.14$. Fair hair and light
eyes co-occur far more often than independence would predict—the variables are
strongly associated.
\end{example}

\section{Expectation of a Function; Covariance and Correlation}

For a function $g(X,Y)$, $\E[g(X,Y)]=\sum_{a,b} g(a,b)\,p(a,b)$. The most
important case is the \emph{covariance}, which measures how $X$ and $Y$ vary
together:
\[
  \Cov(X,Y)=\E[XY]-\E[X]\E[Y].
\]
Its scale-free version is the \emph{correlation coefficient}
\[
  \rho(X,Y)=\frac{\Cov(X,Y)}{\sqrt{\Var(X)\Var(Y)}}\in[-1,1].
\]
Independent variables are uncorrelated ($\rho=0$), but the converse can fail.

\begin{example}[Covariance from a joint table]
For the joint distribution of $U,V$ with
$\E[U]=1$, $\E[V]=\tfrac12$, and $\E[UV]=\tfrac12$, compute $\Cov(U,V)$ and
$\rho(U,V)$.

\solution $\Cov(U,V)=\E[UV]-\E[U]\E[V]=\tfrac12-(1)(\tfrac12)=0$. Since both
variances are positive, $\rho(U,V)=0$ as well. Here zero covariance happens to
coincide with the variables being uncorrelated by construction.
\end{example}

\begin{example}[Uncorrelated but dependent, and variances of sums]
Consider $X\in\{0,1,2\}$, $Y\in\{0,1\}$ with joint distribution
\[
\begin{array}{c|ccc|c}
 & 0 & 1 & 2 & p_Y\\\hline
 Y=0 & 1/6 & 1/6 & 1/6 & 1/2\\
 Y=1 & 1/4 & 0 & 1/4 & 1/2\\\hline
 p_X & 5/12 & 1/6 & 5/12 & 1
\end{array}
\]
Show $X$ and $Y$ are dependent yet uncorrelated, and find $\Var(X+Y)$ and
$\Var(X-Y)$.

\solution One computes $\E[XY]=\tfrac12$, $\E[X]=1$, $\E[Y]=\tfrac12$, so
$\Cov(X,Y)=\tfrac12-(1)(\tfrac12)=0$: uncorrelated. Yet they are dependent, since
$\Prob(X=1,Y=1)=0\neq\Prob(X=1)\Prob(Y=1)$. Thus \emph{zero correlation does not
imply independence}—correlation only detects \emph{linear} association. With
$\Var(X)=\tfrac56$ and $\Var(Y)=\tfrac14$, and using
\[
  \Var(X\pm Y)=\Var(X)+\Var(Y)\pm 2\Cov(X,Y),
\]
both $\Var(X+Y)$ and $\Var(X-Y)$ equal $\tfrac56+\tfrac14=\tfrac{13}{12}$,
because the covariance term vanishes.
\end{example}

\begin{example}[Tuning the correlation]
Let $X,Y\in\{0,1\}$ have joint distribution depending on a parameter $\epsilon$
with $-\tfrac14\le\epsilon\le\tfrac14$:
\[
\begin{array}{c|cc|c}
 & 0 & 1 & p_X\\\hline
 0 & 1/4-\epsilon & 1/4+\epsilon & 1/2\\
 1 & 1/4+\epsilon & 1/4-\epsilon & 1/2\\\hline
 p_Y & 1/2 & 1/2 & 1
\end{array}
\]
Express $\rho(X,Y)$ in terms of $\epsilon$ and find the values giving
$\rho=1,0,-1$.

\solution The margins are uniform on $\{0,1\}$, so $\E[X]=\E[Y]=\tfrac12$ and
$\Var(X)=\Var(Y)=\tfrac14$ regardless of $\epsilon$. Computing
$\E[XY]=\tfrac14-\epsilon$ gives $\Cov(X,Y)=(\tfrac14-\epsilon)-\tfrac14=-\epsilon$,
hence
\[
  \rho(X,Y)=\frac{-\epsilon}{\sqrt{(1/4)(1/4)}}=-4\epsilon.
\]
So $\rho=1$ when $\epsilon=-\tfrac14$ (the off-diagonal probabilities vanish and
$X=Y$ deterministically), $\rho=0$ when $\epsilon=0$ (independence), and
$\rho=-1$ when $\epsilon=\tfrac14$ ($X=1-Y$). The single parameter $\epsilon$
sweeps the correlation across its entire range.
\end{example}

\section{Continuous Joint Distributions}

For continuous $X,Y$ the joint density $f(x,y)$ satisfies
$\Prob((X,Y)\in A)=\iint_A f\,dx\,dy$, the marginals are
$f_X(x)=\int f(x,y)\,dy$ and $f_Y(y)=\int f(x,y)\,dx$, and independence is
$f(x,y)=f_X(x)f_Y(y)$ (equivalently $F(x,y)=F_X(x)F_Y(y)$).

\begin{example}[From a joint cdf]
Suppose $F(x,y)=1-\tfrac1x-\tfrac1{y^2}+\tfrac1{xy^2}$ for $x\ge 1,\,y\ge 1$
(and $0$ otherwise). Find the joint density, the marginal densities, and decide
whether $X$ and $Y$ are independent.

\solution The joint density is the mixed partial derivative,
\[
  f(x,y)=\frac{\partial^2 F}{\partial x\,\partial y}=\frac{2}{x^2 y^3},\qquad x\ge 1,\ y\ge 1.
\]
Integrating out one variable at a time,
$f_X(x)=\int_1^\infty \tfrac{2}{x^2 y^3}\,dy=\tfrac1{x^2}$ and
$f_Y(y)=\int_1^\infty \tfrac{2}{x^2 y^3}\,dx=\tfrac{2}{y^3}$. Since
$f_X(x)f_Y(y)=\tfrac{2}{x^2 y^3}=f(x,y)$, the variables are \emph{independent}.
(Equivalently the cdf factors as $F(x,y)=(1-\tfrac1x)(1-\tfrac1{y^2})$.)
\end{example}

\begin{example}[A product-form density]
Let $f(x,y)=12x(1-x)y$ for $0\le x\le 1,\,0\le y\le 1$. Find the marginal
densities, the joint cdf, and $\Cov(X,Y)$.

\solution Integrating out $y$, then $x$,
\[
  f_X(x)=\int_0^1 12x(1-x)y\,dy=6x(1-x),\qquad
  f_Y(y)=\int_0^1 12x(1-x)y\,dx=2y.
\]
Because $f(x,y)=f_X(x)f_Y(y)$, the variables are independent, and the joint cdf
factors:
\[
  F(x,y)=\Bigl(3x^2-2x^3\Bigr)\,y^2,\qquad 0\le x,y\le 1.
\]
Independence forces $\Cov(X,Y)=0$. We can confirm directly:
$\E[X]=\tfrac12$, $\E[Y]=\tfrac23$, and $\E[XY]=\tfrac13=\E[X]\E[Y]$, so indeed
$\Cov(X,Y)=0$.
\end{example}

\section{Linearity of Expectation in Action}

Expectation is linear even when variables are dependent, and this single fact
solves problems that look forbidding.

\begin{example}[Group testing of blood samples]
The blood of $1000$ people must be screened for a disease that each carries
independently with probability $p=0.001$. Rather than $1000$ individual tests,
pool the samples into $25$ groups of $40$: test each pooled sample, and if it is
positive, retest all $40$ individuals (so that group costs $41$ tests). Let
$X_i$ be the number of tests for group $i$. Find the expected total number of
tests.

\solution For a group, $X_i=1$ if the pool is clean and $X_i=41$ otherwise. The
pool is clean with probability $(1-0.001)^{40}\approx 0.96$, so
$\Prob(X_i=1)\approx 0.96$ and $\Prob(X_i=41)\approx 0.04$. Then
\[
  \E[X_i]=1(0.96)+41(0.04)=2.6,
\]
and by linearity the expected total over all $25$ groups is
$\E\!\left[\sum_{i=1}^{25}X_i\right]=25(2.6)=65$. Pooling replaces $1000$
guaranteed tests with about $65$ expected tests—a striking economy that exploits
the rarity of the disease.
\end{example}
